% !TEX root = ../bb_AiRa_Kappaleet.tex

% ö = \%c3\%b6
% ä = \%C3\%A4

\xdef\sectiontitle{\IVsectiontitle} %aiheuttama
\TOCrefs

%%%%%%%%%%%%%%%
\begin{frame}[label=4_1_a]%label= ; [label=4_0_TOC1]
\frametitle{\texorpdfstring{\culine[\sectiontitleulinecolor]{\sectiontitle}}{\sectiontitle}}
\label{\TOCname}

\vspace{-3mm}

\begin{itemize}[leftmargin=0.5cm,itemsep=\chapterTOCitemsep]
    \item[4.1]<1-> \hyperlink{4_1_a}{\CDAlert[blue]{\IVaSubsectiontitle}}
    \item[4.2]<1-> \hyperlink{4_2_a}{\IVbSubsectiontitle}	
    \item[4.3]<1-> \hyperlink{4_3_a}{\IVcSubsectiontitle}	
    \item[4.4]<1-> \hyperlink{4_4_a}{\IVdSubsectiontitle}
    \item[4.5]<1-> \hyperlink{4_5_a}{\IVeSubsectiontitle}
    \item[4.6]<1-> \hyperlink{4_6_a}{\IVfSubsectiontitle}
    \item[4.7]<1-> \hyperlink{4_7_a}{\IVgSubsectiontitle}
\end{itemize}

%\endofslide 
\end{frame}
%%%%%%%%%%%%%%%

%\xdef\sectiontitle{Prelude to fundamental particles} 
\xdef\sectiontitle{\IVaSubsectiontitle} %aiheuttama

%%%%%%%%%%%%%%%
\begin{frame}
\frametitle{\texorpdfstring{\culine[\sectiontitleulinecolor]{\sectiontitle}}{\sectiontitle} }
\appear{}

\semismall
\begin{itemize}[leftmargin=2cm,labelsep=1cm]%,align=left
	\item[450 eea.] Demokritos: \appear{Pienin särkymätön aineen osanen $\oldRightarrow$ \CDAlert[tunired]{atomi} }
	
	\item[1808] Dalton: \appear{\CDAlert[red]{Atomit} aineiden perustavanlaatuisina rakennuspalikoina}
	
	\item[1879] Thomson: \appear{myös \CDAlert[blue]{elektronit} löydettiin, atomeissa on elektroneita }
	
	\item[1913] Rutherford: \appear{\CDAlert[green]{protonit} löydettiin vetyatomista}
	
	\item[1923] Compton: \appear{Myös \CDAlert[fuchsia]{fotonit}  \Alert[fuchsia]{siroavat} kuin ne olisivat \CDAlert[fuchsia]{hiukkasia}} 
	
	\item[1932] Chadwick: \appear{\CDAlert[blue]{neutronit} löydettiin selittämään esim.\ ytimien mitatut massat }
	
	\item[\Alert{1932$\,\rightarrow$}] \Alert[red]{Positroni}\appear{, \Alert[red]{muoni}}\appear{, \Alert[red]{pioni}, \Alert[red]{sigma}, \Alert[red]{rho}, jne. ...}

\item[1964] \CDAlert[fuchsia]{Gell-Mann} \& Zweig: \appear{ \CDAlert[fuchsia]{Kvarkkien} olemassaolon ennustaminen.} \appear{ Protonien, neutronien jne.\ koostuminen kvarkeista} \appear{(nukleoneilla on sisäinen rakenne)} 
\end{itemize}
\vspace{3mm}

\implies{ Lopputuloksena meillä on \CDAlert[fuchsia!80!black]{nykyään} niin sanottu \CDAlert[fuchsia]{hiukkasfysiikan standardimalli} }

\endofslide 
%\endofsection
\end{frame}
%%%%%%%%%%%%%%%

%%%%%%%%%%%%%%%
\begin{frame}
\frametitle{\texorpdfstring{\culine[\sectiontitleulinecolor]{\sectiontitle}}{\sectiontitle} }
\appear{}

% 6.5 cm = midway, 10 cm = 3/4 
%\vspace{2.5mm}
\begin{minipage}{7.5cm}
\begin{itemize}
	\item Tavallinen aine koostuu \CDAlert[fuchsia]{protoneista}\appear{, \CDAlert[red]{neutroneista}} \appear{ja \CDAlert[blue]{elektroneista}}

\item Viimeisten 70 vuoden aikana\appear{, kokeita on tehty \CDAlert[fuchsia]{hiukkaskiihdyttimillä}} \appear{alati kasvavilla energioilla}\appear{, paljastaen satoja muitakin '\CDAlert[fuchsia]{alkeishiukkasia}'}
\implies{ esim. \CDAlert[fuchsia]{kvarkkien} löytyminen }

\item Alkeishiukkasten ja niiden rakenteen ymmärtäminen on muuttunut huomattavasti vuosien kuluessa 
%Understanding of the fundamental particles and their structure have changed during the years
\end{itemize}
\end{minipage}
%\vspace{2.5mm}


\xdef\loc{south east}
\begin{tikzpicture}[remember picture,overlay] % anchor=south
    \node (A) [yshift=-0.25*\figYshift,xshift=\figXshift, anchor=\loc] at (current page.\loc) {\includegraphics[height=9.5cm]{Figures_booklet/SOM_11_Fundamental_particles_figures/Tikz_SOM_Fundamental_figures.pdf}}; % 8cm max height
\end{tikzpicture}

\endofslide 
%\endofsection
\end{frame}
%%%%%%%%%%%%%%%

%%%%%%%%%%%%%%%
\begin{frame}
\frametitle{\texorpdfstring{\culine[\sectiontitleulinecolor]{\sectiontitle}}{\sectiontitle} }
\framesubtitle{Voiman kantama?}%Range of a force?
\appear{}

\begin{itemize}
	\item \CDAlert[blue]{Alkeishiukkaset vuorovaikuttavat vain muutamalla} \CDAlert[blue]{tyypillisellä tavalla}

\item Mekaniikassa voima tarvitsee usein fyysisen kontaktin välittyäkseen \appear{(jännitys, kitka, ...)}\appear{, mutta gravitaatio vaikuttaa tyhjiön läpi}\appear{, kuten myös sähkömagneettiset voimat!}

\item Nämä \appear{\CDAlert[red]{kenttävoimat}} \appear{ovat pikemminkin sääntö kuin poikkeus} %\appear{Such forces can be called as \CDAlert[red]{field forces}}%\CDAlert[red]{voimakentät}

\item Lisäksi kontaktia vaativat voimat (jännitys, kitka, jne.) \appear{paljastuvat mikroskooppisessa tarkastelussa sähkömag-\\
neettisiksi voimiksi}

%\item More generally, having no contact is more of a rule than an exception. Also electromagnetic force acts through a vacuum. Gravity and electromagnetic forces are so called field forces. Furthermore, contact forces (tension etc.) eventually turn out to be a sum of microscopic electromagnetic forces
\end{itemize}

% 6.5 cm = midway, 10 cm = 3/4 
\vspace{2.5mm}
\begin{minipage}{8.0cm}
\begin{itemize}
	\item \CDAlert[fuchsia]{Vahva voima} vaikutti myös 'etäisyyksien päähän'%
	\appear{\xdef\loc{south east}%
\begin{tikzpicture}[remember picture,overlay]% anchor=south
    \node (A) [yshift=-0.45*\figYshift,xshift=\figXshift, anchor=\loc] at (current page.\loc) {\includegraphics[width=5cm]{Figures_booklet/SOM_11_Fundamental_particles_figures/SOM_Nuclear_Table_2}};% 8cm max height
\end{tikzpicture}\CR[ref=h]{}}%
\appear{, vaikkakin sen \CDAlert[fuchsia]{kantama oli} \CDAlert[fuchsia]{erittäin lyhyt}}
\end{itemize}
\end{minipage}
%\vspace{2.5mm}

\endofslide 
%\endofsection
\end{frame}
%%%%%%%%%%%%%%%



%%%%%How Forces Act?
\xdef\sectiontitle{\IVaSubsectiontitle} 

%%%%%%%%%%%%%%%
\begin{frame}
\frametitle{\texorpdfstring{\culine[\sectiontitleulinecolor]{\sectiontitle}}{\sectiontitle} }
\framesubtitle{Moderni tarkastelutapa}
\appear{}

\begin{itemize}
%	\item \CDAlert[blue]{Moderni näkökulma}:
%\item[] \begin{itemize}
	 \item \CDAlert[tuniblue]{Voimat vaikuttavat vaihtamalla välittäjähiukkasia}
\item Jokaisella voimalla \appear{(4:stä perusvoimasta)} \appear{on oma välittäjähiukkanen}
\item Välittäjähiukkanen tunnetaan myös kyseisen kvanttikentän kvanttina
%\end{itemize}
\implies{  Kvanttikenttäteoria (QFT) }

\item \CDAlert[red]{Kuinka vuorovaikutus selitetään?}

\item On olemassa kahdenlaisia hiukkasia:
\item[] \begin{itemize}
	\item Hiukkaset jotka \CDAlert[blue]{kokevat} voiman
\item Hiukkaset joiden vaihtaminen ilmentää voimaa\appear{, eli vuorovaikutuksen \CDAlert[tuniblue]{välittäjähiukkaset}} %Those particles which are exchanged to convey the force\appear{, i.e. to \CDAlert[tuniblue]{mediate} the interaction}

\end{itemize}

\item Huomaa että joissain tilanteissa välittäjähiukkaset kokevat itsekin välittämänsä voiman!
\end{itemize}

\endofslide 
%\endofsection
\end{frame}
%%%%%%%%%%%%%%%


%%%%%%%%%%%%%%%
\begin{frame}
\frametitle{\texorpdfstring{\culine[\sectiontitleulinecolor]{\sectiontitle}}{\sectiontitle} }
\appear{}

\begin{minipage}{7.5cm}
\begin{itemize}
	\item Kaksi fysiikanopiskelijaa\appear{, Anna ja Bob}\appear{, jotka liukuvat kitkattomalla jäällä}%
	\appear{\xdef\loc{east}%
\begin{tikzpicture}[remember picture,overlay] % anchor=south
    \node (A) [yshift=-0*\figYshift,xshift=\figXshift, anchor=\loc] at (current page.\loc) {\includegraphics[height=8cm]{Figures_booklet/SOM_11_Fundamental_particles_figures/SOM_Fundamental_particles_Figure_1.png}};% 8cm max height
\end{tikzpicture}\CR[]{}%
}%
\appear{, vuorovaikuttavat niin että Anna heittää lumipallon Bobille, joka ottaa sen kiinni}

\item Molempien henkilöiden liikemäärä muuttuu: \appear{Annan liikemäärä muuttuu $-\vec{p}_{\text {pallo}}$}\appear{, ja Bobin liikemäärä muuttuu  $+\vec{p}_{\text {pallo}}$}

\item Tämä näyttää "\Alert[fuchsia]{lumipallon välittämältä hylkivältä voimalta}"\appear{, missä pallo siirtää Annan liikemäärää Bobille}

\item \CDAlert[blue]{Feynmanin kaavio/diagrammi} kuvastaa tätä vuorovaikutusta
\end{itemize}
\end{minipage}

\endofslide 
%\endofsection
\end{frame}
%%%%%%%%%%%%%%%



%%%%%%%%%%%%%%%
\begin{frame}
\frametitle{\texorpdfstring{\culine[\sectiontitleulinecolor]{\sectiontitle}}{\sectiontitle} }
\appear{}

%\seminormal
\begin{itemize}[itemsep=1.7mm]
	\item Esimerkki matkii onnistuneesti Annan ja Bobin välistä hylkivää voimaa\appear{, mutta \CDAlert[fuchsia]{analogia ei selitä puoleensavetäviä voimia}}
	\appear{\xdef\loc{south east}
\begin{tikzpicture}[remember picture,overlay] % anchor=south
    \node (A) [yshift=-0.35*\figYshift,xshift=\figXshift, anchor=\loc] at (current page.\loc) {\includegraphics[scale=0.85,page=2]{Figures_booklet/SOM_11_Fundamental_particles_figures/Tikz_SOM_Fundamental_figures_combo.pdf}}; % 8cm max height
\end{tikzpicture}}
	\item Kuvatakseen puoleensavetävää voimaa\appear{, lumipallon heiton jälkeen}\appear{, Annan pitäisi liikkua Bobia kohti}\appear{, ja Bobin ottaessa pallon kiinni}\appear{, hänen täytyisi liikkua Annaa kohti} 
	\implies{ \CDAlert[fuchsia]{Liikemäärä ei säilyisi!}}

\item Esimerkissä\appear{ Annalla oli lumipallo jo ennen heittoa}\appear{, ja kun Bob otti pallon kiinni}\appear{, hän säilytti sen}
\end{itemize}

% 6.5 cm = midway, 10 cm = 3/4 
\vspace{1.7mm}
\begin{minipage}{9.0cm}
\begin{itemize}[itemsep=1.7mm]
	\item Todellisuudessa lumipallo \appear{(välittäjähiukkasena)} \appear{\CDAlert[red]{syntyisi vain} \CDAlert[red]{lennon ajaksi}} \appear{ja se \CDAlert[tunired]{annihiloituisi} välittömästi \CDAlert[tunired]{kiinnioton jälkeen}}
	\implies{\CDAlert[blue]{Energia ei näytä säilyvän tällaisessa}\\ \CDAlert[blue]{makroskooppisessa esimerkissä!}}
\end{itemize}
\end{minipage}
%\vspace{2.5mm}

\endofslide 
%\endofsection
\end{frame}
%%%%%%%%%%%%%%%

%%%%%%%%%%%%%%%
\begin{frame}
\frametitle{\texorpdfstring{\culine[\sectiontitleulinecolor]{\sectiontitle}}{\sectiontitle} }
\framesubtitle{Virtuaalihiukkaset}
\appear{}

\semismall
\begin{itemize}[itemsep=1.5mm]
	\item Välittäjähiukkaset ovat \CDAlert[fuchsia]{virtuaalihiukkasia}\appear{, ne ovat olemassa vain vuorovaikutuksen ajan}\appear{, eikä niiden tarvitse toteuttaa tavallisia energian ja liikemäärän säilymislakeja}\appear{, kuten 'oikeiden' hiukkasten tarvitsee}
	
	\item Ilmiö muistuttaa tunneloitumista\appear{, %the kinetic energy of particle is negative, and 
	missä virtuaalihiukkasten esiintymistodennäköisyys pienenee} \appear{kun oikeiden hiukkasten välinen etäisyys kasvaa}
\end{itemize}

% 6.5 cm = midway, 10 cm = 3/4 
\vspace{1.5mm}
\begin{minipage}{8.3cm}
\begin{itemize}[itemsep=1.5mm]
%	\item[] The phenomenon resembles tunneling, the kinetic energy of particle is negative, and the probability for the occurence of the particle decreases when the separation between the interacting particles increases

%\item This is the mechanism which \CDAlert[red]{H. Yukawa} (1935) suggested on how the strong nuclear force would act between two nucleons, the nucleons exchange virtual particles
\item Tätä \CDAlert[fuchsia]{virtuaalihiukkasten vaihtoa} \appear{ehdotettiin kahden nukleonin välisen ydinvoiman mekanismiksi}\appear{ vuonna 1935 (\CDAlert[red]{H. Yukawa}) }
\appear{\xdef\loc{south east}%
\begin{tikzpicture}[remember picture,overlay] % anchor=south
    \node (A) [yshift=-0.35*\figYshift,xshift=\figXshift, anchor=\loc] at (current page.\loc) {\includegraphics[scale=0.85,page=1]{Figures_booklet/SOM_11_Fundamental_particles_figures/Tikz_SOM_Fundamental_figures_combo.pdf}};% 8cm max height
\end{tikzpicture}}%

\item Vahvan voiman kantama on \CDAlert[red]{lyhyt} \appear{koska  \CDAlert[red]{välittäjähiukkasella on massa}.} \appear{Ajatusta voidaan soveltaa myös sähkömagneettiseen vuorovaikutukseen.} \appear{Sähköiset varaukset lähettävät jatkuvasti ympärilleen virtuaalihiukkasia}
 
\item Kantama on äärettömän pitkä koska virtuaalihiukkaset ovat \CDAlert[blue]{fotoneita}\appear{, joilla \CDAlert[blue]{ei ole} \CDAlert[blue]{lepomassaa}}
\end{itemize}
\end{minipage}
%\vspace{2.5mm}

\endofslide 
%\endofsection
\end{frame}
%%%%%%%%%%%%%%%

%%%%%%%%%%%%%%%
\begin{frame}
\frametitle{\texorpdfstring{\culine[\sectiontitleulinecolor]{\sectiontitle}}{\sectiontitle} }
\framesubtitle{Virtuaaliset välittäjähiukkaset}
\appear{}

%Mediating particle \item mass vs. range
\begin{itemize}
	\item Mekanismin ja \CDAlert[tuniblue]{virtuaaalisten välittäjähiukksten} olemassaolo \appear{voidaan ymmärtää soveltamalla Heisenbergin \CDAlert[blue]{epätarkkuusperiaatetta} aika–energia (konjugaatti)muuttujapariin}

\item Oletetaan että virtuaalihiukkasilla on tietty aikaväli $\Delta t$ vuorovaikuttaa \appear{(eli ajan epätarkkuus)}\appear{, silloin hiukkasen energiaan $\Delta E$ liittyvän epätarkkuuden täytyy toteuttaa yhtälö } $\appear{\Delta t \Delta E} \approx \appear{\hbar}$

\item Hiukkasen energia voi vaihdella välillä $\Delta E$ säilymislakien rikkoontumatta. \appear{Toisin sanoen jopa energia $\Delta E=m c^{2}$ on sallittu virtuaalihiukkasille}\appear{, jos sillä on lepomassa}

\item Tällöin $\Delta t \Approx \hbar / m c^{2}$\appear{, ja mahdolliseksi kantamaksi saadaan} $\appear{\Delta x} \approx \appear{c \Delta t} \approx \appear{\hbar / m c}$ 

\item Yhteenvetona\appear{ \CDAlert[fuchsia]{virtuaalihiukkasten kantama} $\Approx \Frac{\hbar}{c} \Frac{1}{m}$}
\end{itemize}

\endofslide 
%\endofsection
\end{frame}
%%%%%%%%%%%%%%%

%%%%%%%%%%%%%%%
\begin{frame}
\frametitle{\texorpdfstring{\culine[\sectiontitleulinecolor]{\sectiontitle}}{\sectiontitle} }
\framesubtitle{Välittäjähiukkanen, massa vs. kantama}%Virtual mediating particles, mass vs. range
\appear{}

\seminormal
\begin{itemize}
	\item \CDAlert[red]{Mitä suurempi massa} välittäjähiukkasella on\appear{, sitä \CDAlert[red]{lyhyempi kantama} on vuorovaikutuksella}\appear{, ja päinvastoin} 
	
	\item Vahvan vuorovaikutuksen kantama on ${\sim\;}1 \mathrm{\;fm}$ 
	\implies{ Tästä saadaan arvio mesonin massalle:
\[
m  \equal \appear{3,5 \times 10^{-28} \mathrm{\;kg}} \approx \appear{300 \mathrm{m}_{\mathrm{e}}} \approx \appear{200 \mathrm{\;MeV} / \mathrm{c}^{2}}
\] }
\item Olettamalla välittäjähiukkanen fotoniksi\appear{, kuten sähkömagneettisissa vuorovaikutuksissa, kantamaksi saadaan:}
\[
\text {kantama} \approx \appear{c} \appear{\Delta t} \approx \frac{c \hbar}{\Delta E}  \equal \appear{c \hbar} \frac{\lambda}{h c} \equal \frac{\lambda}{2 \pi}
\]

\item Koska fotoneilla ei ole (lepo)massaa\appear{, niiden energioille ei ole alarajaa}\\
\appear{(}$\Rightarrow$ \appear{kantama voi olla erittäin pitkä, lähestyen äärettömyyttä}\appear{)}

\item Itse asiassa kahden elektronin \CDAlert[blue]{vuorovaikuttaessa sähköisesti}\appear{, molemmat elektronit lähettävät jatkuvasti \CDAlert[blue]{virtuaalisia fotoneita} joita toiset elektronit absorboivat}
\end{itemize}

%\endofslide 
\endofsection
\end{frame}
%%%%%%%%%%%%%%%
